For fixed \(S\) and arm \(\rho\), membership in \(\operatorname{PO}(S)\) is exactly feasibility of
\[
q\ge0,\quad A_Dq=m_D,\quad R_eq=p_e\ (e\in S),\quad
(n_{\rho,D}-n_{\sigma,D})^{\mathsf T}q\ge0\ (\sigma\in\mathcal R_{\mathrm{arm}}).
\tag{1}
\]
The common denominator is certified positive, so (1) is equivalent to value comparison.  Finite subset enumeration and exact rational or real-algebraic linear feasibility compute every possibly-optimal set and every minimum-cost solution.

For the obstruction, take the binary DAG \(T_{\mathrm{tr}}\to Y\), intervene on \(T_{\mathrm{tr}}\), and use the four deterministic response maps
\[
f_{00}(t)=0,\quad f_{01}(t)=t,\quad f_{10}(t)=1-t,\quad f_{11}(t)=1,
\]
ordered as \(\omega_1=(0,0),\omega_2=(0,1),\omega_3=(1,0),\omega_4=(1,1)\).  Let sure evidence be common, let
\[
B_\rho=\{(Y_0,Y_1)\ne(1,1)\},\qquad
B_\sigma=\{(Y_0,Y_1)=(1,1)\}.
\]
Declare two typed partial-copy descriptors \(a_1,a_2\).  Experiment \(e_1\) reveals their genuine two-world cell \(C_1=\{Y_0=Y_1\}\), and experiment \(e_2\) reveals the genuine two-world cell \(C_2=\{(Y_0,Y_1)\ne(0,0)\}\).  Neither cell indicator is a function of \(Y_0\) alone or of \(Y_1\) alone.  Exhaustive evaluation on all four response maps gives
\[
\begin{array}{c|rrrr}
&n_\rho&n_\sigma&r_1&r_2\\\hline
\omega_1&1&0&1&0\\
\omega_2&1&0&0&1\\
\omega_3&1&0&0&1\\
\omega_4&0&1&1&1
\end{array}
\tag{2}
\]
so the rows are derived from concrete binary counterfactual events.  Use normalization as the only base equation and let the two typed experiment cells have supplied values \(P_{a_1}(C_1)=P_{a_2}(C_2)=1\), each at unit cost.  The joint menu is realized by the original-alphabet response-block law \(q=\delta_{\omega_4}\), and
\[
\mathcal K_{\{e_1\}}=\operatorname{conv}\{\delta_1,\delta_4\},\quad
\mathcal K_{\{e_2\}}=\operatorname{conv}\{\delta_2,\delta_3,\delta_4\},\quad
\mathcal K_{\{e_1,e_2\}}=\{\delta_4\}.
\]
At \(\delta_1,\delta_2,\delta_3\), \(\rho\) is uniquely optimal; at \(\delta_4\), \(\sigma\) is uniquely optimal.  Thus each singleton experiment retains both arms, whereas their intersection retains only \(\sigma\).  This proves the four possibly-optimal sets and the increasing marginal gain, hence non-submodularity.  Both experiments are necessary, giving \(x_1+x_2\ge2\); the feasible binary point \((1,1)\) has cost two, and multiplier one proves the same lower bound for the relaxation.